% =====================================================================================================
% Section II. What Optimizing Labs' Data Identify  (version 3, after referee round 2)
% Writer: section "identification" (writer3, 2026-09-24). v2 text: sections/v2/identification.tex.
% Binding sources: paper/notes/revision_plan_v3.md (Section 2: four formal results; Prop 2 with the rate as a Remark;
% Prop 3 = model-free sigma* and w + the beyond-design sign result; notation varsigma, e_p, no b = 1 - a),
% paper/notes/theory_main_text_v2.md, appendix_proofs.tex (Props. A1, A2, A10, A11; Remark A-grid),
% output/memos/m6_montecarlo_v2.md (Figure 2), rb1_sigmaC.md and rb2_decisions.md (+ reviews).
% Labels: prop:ident (Prop 2), rem:rate (Remark 1), prop:modelfree (Prop 3), prop:pi = alias of Prop 3(v).
% Numbers are verified against output/tables/*.csv and data/processed/rb2_decisions/*.csv; ledger in
% paper/notes/writer3_identification.md. Figure 2 is m6_montecarlo_fig2_v2 re-rendered with the v3 notation
% (varsigma) by code/paper/make_ident_v3.py; same data.
% Round 3 (WP3, fix list P4-P8, P12-P13): Prop. 3(v) with the interval anchor, joint bands and union of paths; Remark 1
% scope; Prop. 2 wording; Figure 2 notes; general-isocost sentence; numbers from output/tables/rb5_sign_paper_numbers.csv.
% [round-3 integration, 2026-09-25] Word-cap cuts in the text only (no number, formal result or requested item removed);
% proof pointer added to the opener. Log: paper/notes/round3_integration_log.md.
% =====================================================================================================

\providecommand{\E}{\mathbb{E}}
% \identpartlabel{(v)}{key}: a label that refers to one part of the current proposition, so that \ref{key} prints
% "3(v)". Used to keep the v2 key prop:pi (partial identification beyond the design) pointing at Proposition 3(v).
\makeatletter
\providecommand{\identpartlabel}[2]{\protected@edef\@currentlabel{\theproposition#1}\label{#2}}
\makeatother

\section{What Optimizing Labs' Data Identify}\label{sec:ident}

Reading developers' choices requires $\sigma^*$ and $M^*(C)$, the scale and the zero point of the wedge in
equation~\eqref{eq:wM}. This section asks which data identify them; proofs
are in Online Appendix~\ref{app:proofs}.

\subsection{Optimization Removes the Variation That Identifies Curvature}\label{sec:ident-onpath}

A sample is \emph{on-path} if its runs share one technology, each run minimizes training compute for its loss (every
wedge equals one), and compute varies over an interval.

\begin{proposition}[What optimizing labs' data identify]\label{prop:ident}
Let runs follow \eqref{eq:tech} with $\psi_N=\psi_D=0$.
\begin{enumerate}
\item[(i)] (On the path.) In an on-path sample, $n_i$ and $d_i$ are affine in $c_i$, with slopes $a$ and $1-a$:
choices reveal the path, that is, $a$ and $M^*(C)$, and nothing else. Outcomes satisfy
$\E[y\mid c]=\gamma c+\mathrm{const.}+\E[\omega\mid c]$, so within the model $\gamma$ is identified if and only if
$\E[\omega\mid c]$ is known up to a constant, as in a single-lab sweep. Curvature is not identified: for every
$\sigma^*\in(0,1)$ some member of family \eqref{eq:tech} generates the same joint distribution of $(n,d,y)$, and by
Lemma~\ref{lemma:sigma} no value $\sigma^*>1$ does (in family \eqref{eq:tech}, none with $\sigma^*\ge1$).
\item[(ii)] (Global versus local identification.) Suppose also that $\E[\omega\mid c]$ is constant and $\kappa=1$.
Globally, with $A,B>0$, $\alpha=\gamma/a$ and $\beta=\gamma/(1-a)$ are identified by the functional form, and $A/B$, hence the
technology's own cost-minimizing mix, only by imposing that the observed path is optimal. Locally, identification fails:
every perturbation that moves the decay rates of the two loss terms along the path, $\alpha a$ and $\beta(1-a)$, in
opposite directions leaves on-path outcomes unchanged to first order once the common level is re-split between the two
terms (at a given split only one such direction is flat). The least-squares criterion, profiled over that split, is
therefore of fourth order along such directions, and $\sigma^*$ is not identified at first order, whether or not $\alpha=\beta$. Imposing
$\beta/(\alpha+\beta)=a$ restores first-order identification.
\item[(iii)] (Off the path.) Let runs deviate from the path with log wedges $\ln w_i$, with $E$ known and Gaussian noise.
To leading order, the information about the path's location ($\ln G$, and so $\ln M^*$ at a given compute) is
proportional to the residual dispersion of the $\ln w_i$. The
information about $\sigma^*$ is bounded by a multiple of $\sum_i\Phi_i^2$, where $\Phi_i\equiv\ln\big(C_i/C_{\min}(L_i)\big)$
is run $i$'s Farrell allocative loss: the log of its compute relative to the least compute that attains its loss. If the
deviations of a fixed design are scaled by $\varsigma\to0$, the two informations are $O(\varsigma^2)$ and
$O(\varsigma^4)$, so the standard error of $\ln\hat M^*$ grows at least like $\varsigma^{-1}$ and that of $\hat\sigma^*$ at
least like $\varsigma^{-2}$, exactly so for designs that cross several compute levels with symmetric transverse
offsets.
\end{enumerate}
\end{proposition}

A ray has no curvature. Facing the same log-cost weights, cost-minimizing labs that share a technology place every run
on the expansion path, where reducible loss is log-linear in compute. Choices trace the path
and outcomes the frontier, but neither reveals how fast loss rises off the path, which is what $\sigma^*$ measures
(Figure~\ref{fig:geometry}, panel~b): at five times the compute-optimal ratio, reducible loss lies 2.0, 4.1 and 7.6
percent above the frontier for three technologies that share it and the path, with $\sigma^*=0.85$, 0.74 and 0.60.



\begin{remark}[The rate on the path]\label{rem:rate}
Part (ii) makes the on-path model the analogue of an over-fitted mixture, whose parameters can be learned at best at the
rate $n^{-1/4}$ \citep{chen1995optimal}. We conjecture the same bound here for estimators of $\sigma^*$ from on-path data
that leave $\beta/(\alpha+\beta)$ free, such as unrestricted least-squares fits of the Chinchilla form, whose Wald
intervals, built on the near-singular Hessian, are then invalid. Imposing $\beta/(\alpha+\beta)=a$, which the optimality that defines an on-path sample
implies under $\kappa=1$, removes the singularity: $\sigma^*=2/\{2+\gamma/[a(1-a)]\}$ is then estimated at the usual
root-$n$ rate. The argument is heuristic: the results of
\citet{rotnitzky2000likelihood} for singular information do not apply when $A/B$ is unidentified at the truth
\citep[see also][]{sargan1983identification}. In simulations with 40 on-path compute levels (holding the design fixed and letting the noise
fall, which is equivalent to replicating it $n$ times; Online Appendix Remark~\ref{rem:A-rate}), the 75th percentile of
the error of unrestricted least-squares estimates of $\sigma^*$ shrinks as $n^{-0.24}$, and that of restricted ones as
$n^{-0.50}$.
\end{remark}

Real ladders and labs deviate from the path. Because the frontier is a power law, output decomposes exactly as
\begin{equation}\label{eq:ident-decomp}
y_i \;=\; \gamma c_i+\mathrm{const.}-\gamma\,\Phi_i-\upsilon_i ,
\end{equation}
with $\omega$ normalized to zero. Curvature enters output only through the Farrell losses, which are Harberger triangles,
$\Phi_i\approx\sigma^*(\ln w_i)^2/[4(1-\sigma^*)]$. A run reveals where the path lies to the extent that it is off the
path, first order in its log wedge, and the path's curvature only to the extent that it wastes compute, second order. Information squares these: halving the dispersion of log wedges quarters the information about $M^*$ and
divides the information about $\sigma^*$ by sixteen (the orders are unchanged when $E$ is estimated; Online
Appendix~\ref{app:proofs}).

Proposition~\ref{prop:ident} is the point of \citet{marschak1944random} in its cleanest form: identification requires
the variation in relative prices, adjustment frictions or
optimization errors of \citet{bond2005adjustment} \citep[see also][]{gandhi2020identification}; here it
comes from heterogeneous wedges (Proposition~\ref{prop:wedge}) and belief errors such as allocation by the
\citet{kaplan2020scaling} law. The rates in part (iii) are the standard design rates for a curvature
parameter, whose information grows with the fourth power of the design's spread (\citealt{kiefer1959optimum} on the
highest coefficient of a polynomial regression; \citealt{box1959design} on nonlinear models); what is specific here is
that economic wedges generate the design's spread.

The Chinchilla data show part (i) directly. On the 41 runs within 0.15 log points of the fitted compute-optimal size, the
profile likelihood for $\sigma^*$ with $\kappa$ free is flat: its 95 percent confidence set is $[0.20,0.99]$, the edge of
the grid searched, with $\chi^2_1$ critical values (about $[0.15,0.99]$ with bootstrap ones), against $[0.69,0.71]$
($[0.68,0.72]$) on the full 240-run design, which adds transverse variation (Online Appendix~\ref{app:additional}).

\begin{figure}[tp]
\centering
\includegraphics[width=\textwidth]{m6_montecarlo_fig2_v3}
\caption{Identification by Design: Monte Carlo Evidence}\label{fig:designs}
\begin{figurenotes}[Notes]
Simulations of the \citet{besiroglu2024chinchilla} technology ($\sigma^*=0.737$; $M^*=18.1$ at $10^{24}$ FLOP) with normal
noise of standard deviation 0.0075 in $\ln L$ (the Besiroglu fit's residual s.d.\ is 0.0076). Every design has 90 runs
and the same total compute, $5.1\times10^{22}$ FLOP; except in the factorial design, the runs sit ten on each of nine
budgets matching Chinchilla's nominal IsoFLOP budgets ($6\times10^{18}$ to $3\times10^{21}$ FLOP). Optimizing
labs set $\ln(D/N)$ at its compute-optimal value plus a normal allocation error with standard deviation $\varsigma$; the
IsoFLOP design places ten sizes log-spaced over $[N^*/16,16N^*]$ at each budget, and the factorial design crosses nine
values of $N$ with ten of $D$, each spanning a factor of 256 around the expansion path, so that its runs range from
$1.8\times10^{17}$ to $1.2\times10^{22}$ FLOP. Panel A: root mean squared error of $\hat\sigma^*$ and of $\ln\hat M^*(10^{24})$ for
optimizing labs, 500 replications per design. Primal: fit of the \citet{hoffmann2022training} form with $\kappa=1$ by the
Huber loss with threshold $10^{-3}$, started from nine values drawn at random within the region of Hoffmann et al.'s
initialization grid (exponents on $[0.1,1.2]$ rather than $[0,2]$), none at the truth. Dual: exponents from the path of the optima and the frontier under $\kappa=1$ and
optimality. Horizontal lines: primal estimator under the IsoFLOP and factorial designs. Panel B: median profile
likelihood-ratio statistic for $\sigma^*$ with $\kappa$ free, relative to the unrestricted optimum, on a 29-point grid
over $[0.05,0.99]$ (150 replications), for on-path labs ($\varsigma=0.02$), optimizing labs with $\varsigma=0.3$ and the
IsoFLOP design; horizontal line, the 95 percent $\chi^2_1$ critical value; dotted line, the truth.
\end{figurenotes}
\begin{figurenotes}[Source]
Author's simulations; design, estimators and full tables in Online Appendix~\ref{app:mc}.
\end{figurenotes}
\end{figure}

Figure~\ref{fig:designs} shows what the rates mean at realistic noise and design sizes, in simulated designs of 90 runs
with equal total compute that differ in where the runs sit. With $\varsigma=0$, a Huber fit of the Chinchilla form ends on a corner, with one power-law term switched off and its
exponent not identified at all, in 87 percent of replications (56 percent with 64 random starts), each fitting the data at least as well as the truth.
As $\varsigma$ rises from 0.05 to 0.2, the root mean squared error of $\hat\sigma^*$ falls from 0.089 to 0.015, against
0.0036 for an IsoFLOP design with the same compute (panel A); for $\ln\hat M^*$ it is 4.6 at $\varsigma=0.1$ and 0.59 at
$\varsigma=1$, against 0.16.

With $\kappa$ free (panel B), this precision proves to be functional-form information, as part (ii) predicts: the 95
percent profile confidence set for $\sigma^*$ covers the entire grid searched, $[0.05,0.99]$, in 91 percent of on-path
replications, and at $\varsigma=0.3$ it excludes low elasticities but still contains
every value above 0.9 in 76 percent of replications, whereas the IsoFLOP and factorial designs give sets 0.016 and 0.019
wide. On the path, standard inference also fails, and heteroskedastic noise correlated within budgets or training trunks
leaves the ranking of the designs unchanged (Online Appendix~\ref{app:mc}).

\subsection{Designed Variation Identifies Curvature, and What Survives beyond the Design}\label{sec:ident-modelfree}

The next result reads the technology off movement across the path without a functional form, and says what survives
where no design reaches. Let
$x\equiv\ln M$, so that an IsoFLOP profile is loss as a function of $x$ at fixed $c$, with minimum at
$x^*(c)\equiv\ln M^*(C)$; subscripts denote partial derivatives. A design budget is \emph{bracketed} if at least two of its IsoFLOP
runs lie on each side of the profile's fitted minimum, within the estimation window (Online
Appendix~\ref{app:add-modelfree}).

\begin{proposition}[Model-free $\sigma^*$ and $w$, and the sign beyond the design]\label{prop:modelfree}
Let $L$ be any measure of output that falls as quality rises (loss, bits per byte, or $\ln(L-E)$ for any $E$), and let
the technology be twice continuously differentiable with positive marginal products.
\begin{enumerate}
\item[(i)] (Wedge.) At any point, $w=(L_c-L_x)/(L_c+L_x)$.
\item[(ii)] (Curvature.) At a compute-optimal point,
\begin{equation}\label{eq:ident-mf}
\frac1{\sigma^*}-1 \;=\; \frac{L_{nn}\big|_C}{2\,\lvert dL^*/dc\rvert},
\end{equation}
where $L_{nn}|_C$ is the curvature of the IsoFLOP profile in $\ln N$ at its minimum and $L^*(c)$ is the loss--compute
frontier. The ratio is invariant to strictly monotone transformations of output.
\item[(iii)] (Slope of the log wedge.) At $M^*(C)$, $\partial\ln w/\partial\ln M|_C=1/\sigma^*-1$. In family
\eqref{eq:tech}, $\ln w$ is exactly linear in $\ln(M/M^*(C))$ with this slope; in general the slope varies with $M$.
\item[(iv)] (What $\kappa=1$ imposes.) Under $\kappa=1$ the curvature of every IsoFLOP profile is pinned by the slopes
of the frontier and the path, $-y_{nn}|_C=\gamma^2/[a(1-a)]$; with $\kappa$ free it is a separate parameter.
\item[(v)]\identpartlabel{(v)}{prop:pi} (Beyond the design.) Let $c_J$ be the log of the largest bracketed design budget,
and consider a model with log compute $c>c_J$ and log ratio $x$. Suppose that (a) every IsoFLOP profile of the
developer's technology is single-peaked in $x$ (its only critical point is the peak), which $0<\sigma<1$ everywhere
guarantees (Lemma~\ref{lemma:sigma}); (b)
the anchor $x^*(c_J)$ lies in $[\underline x_J,\overline x_J]$; (c) the path slope $e_p\equiv d\ln M^*/d\ln C$ lies in
$[e_L,e_U]$ between $c_J$ and $c$; and (d) at every compute, the developer's own compute-optimal ratio lies within a
factor $\tau\ge1$ of the anchored technology's. Then the wedge in the developer's first-order conditions satisfies
$\operatorname{sign}(w-1)=\operatorname{sign}\big(x-x^*_{\mathrm{dev}}(c)\big)$, where the developer's compute-optimal
log ratio $x^*_{\mathrm{dev}}(c)$ lies in
\[
\mathcal X^*(c)\equiv\big[\,\underline x_J+e_L(c-c_J)-\ln\tau,\ \ \overline x_J+e_U(c-c_J)+\ln\tau\,\big].
\]
Hence $w>1$ is identified if $x>\sup\mathcal X^*(c)$ and $w<1$ if $x<\inf\mathcal X^*(c)$; when the anchor is a point
and the anchored technology is in family \eqref{eq:tech}, these conditions are also necessary, and otherwise the sign is
not identified. Neither $\sigma^*$ nor a functional form is needed.
\end{enumerate}
\end{proposition}

Part (i) is a change of coordinates: slopes along the isocost and along a ray of constant $M$ give the two marginal
products, whose ratio is the wedge. Part (ii) writes $\sigma^*$ in observables: at the minimum of
an IsoFLOP profile the marginal products of parameters and tokens are equal, the profile's curvature measures how fast
they diverge as the mix moves, and the frontier slope, the marginal
product of compute by the envelope theorem, measures their common level. Because both terms are in the same output
units, the irreducible loss, the transformation of output and $\kappa$ cancel. Both are familiar: the quadratic
coefficient of the parabola fitted to each profile, as in Approach~2 of \citet{hoffmann2022training}, and the slope of
the frontier through the minima; quadratics fitted over wide grids overstate the curvature slightly, which
narrow windows avoid (Online Appendix Remark~\ref{rem:A-grid}; on the corresponding shift of the fitted minimum, see
\citealp{czech2026problems}). With a known cost that is not multiplicative, the same curvature and frontier slope,
weighted by the cost's elasticities, identify $1/\sigma^*$ less the inverse elasticity of substitution of the isocost
itself (Online Appendix Corollary~\ref{cor:A-general}), so the estimator applies to any design that varies two inputs
along known isocosts. Part (iii) gives a test of the family.

Parts (i) to (iv) hold only inside a design's support, but nearly all released models used more compute than
$3\times10^{21}$ FLOP, the largest bracketed budget of any public design (Section~\ref{sec:wedge}); part (v) says what
extrapolating the zero point delivers.

The sign result needs only the location of the minimum: on a single-peaked profile, a model with more tokens per
parameter than is optimal at its compute has a higher marginal product of parameters than of tokens, so $w>1$ whatever
the curvature. Assumption (a) extends the U-shaped profiles of
Lemma~\ref{lemma:sigma} and every bracketing design. The anchor in (b) is a compute-optimal ratio at a bracketed budget,
for Llama~3 at $10^{21}$ FLOP rather than at its unbracketed $3\times10^{21}$ or $10^{22}$, and the path-slope range in
(c) is the substantive extrapolation; Section~\ref{sec:wedge-pi} takes both, with their sampling error, from the designs'
own paths. Normal inputs alone, factor demands that do not fall with compute, imply only
$|e_p|\le1$. The allowance $\tau$ in (d) admits
cross-lab differences in technology (Section~\ref{sec:tech}): a factor-augmenting bias $\chi$ moves $\ln M^*$ by
$2\chi/(\alpha+\beta)$, and the largest tilt across DataDecide's recipes, 0.26,
corresponds to $\tau\approx1.8$ under the reference exponents (we use 1.84) and to 3.4 under DataDecide's own.

Given (b) to (d), a point anchor and an anchored technology in family \eqref{eq:tech}, the set is sharp: every point
in it is the compute-optimal ratio of a technology with positive marginal products and single-peaked profiles that
reproduces the designs up to $c_J$, shifted by a constant factor-augmenting tilt of at most $\ln\tau$ (Online Appendix
Proposition~\ref{prop:A-pi}(ii) and (vii); \citealp{manski2003partial}). With sampling error, a confidence set replaces
the anchor and the slope range; if it is a joint band for the level and slope of a linear path, the set covers
$x^*_{\mathrm{dev}}(c)$ with at least the band's coverage probability, and with several anchored paths the union does
so if the developer's technology lies within $\tau$ of at least one of them (Proposition~\ref{prop:A-pi}(viii)).

What the data identify at frontier scale is therefore a sign, not a level. SmolLM2 1.7B, trained on about 6,400 tokens
per parameter with $1.1\times10^{23}$ FLOP, has a ratio about 47 times the upper end of the set under the anchors of
Section~\ref{sec:wedge}, so its over-training is identified for any tilt allowance below that factor. Llama~3.1 405B, at
38 tokens per parameter and $3.8\times10^{25}$ FLOP, lies inside the set: even the sign of its wedge is not identified.
Magnitudes need the curvature as well: if $\partial\ln w/\partial x|_c\in[\theta_L,\theta_U]$ with $\theta_L>0$, then
for $x>\sup\mathcal X^*(c)$, $\ln w\in\big[\theta_L\big(x-\sup\mathcal X^*(c)\big),\ \theta_U\big(x-\inf\mathcal X^*(c)\big)\big]$
(Online Appendix Proposition~\ref{prop:A-pi}(iv)); in family \eqref{eq:tech} a range for $\sigma^*$ at the model's own
compute gives $[\theta_L,\theta_U]$, and that range is itself an extrapolation.

\textit{Where the results bind.}\label{sec:ident-scope} Proposition~\ref{prop:ident} binds for compute-optimal ladders
and for releases of the compute-optimal era; by the same mechanism, fixed tokens-per-parameter designs run by one lab are
ill-conditioned when $\alpha\approx\beta$ \citep{kricheli2026tokens}, as under the reference technology (DataDecide trains every size to about 100
tokens per parameter; \citealp{magnusson2025datadecide}). Released models since 2023 lie far off
the path, but using their values of compactness as transverse variation requires that these values move the mix without
moving the technology, which factor-biased technology would violate. We therefore take the technology from designed
sweeps and read objectives from released models' wedges.
